\documentclass[12pt]{article}
\usepackage[T1]{fontenc}
\usepackage[utf8]{inputenc}
\usepackage{lmodern}
\usepackage{textcomp}
\usepackage{enumitem}
\usepackage{geometry}
\usepackage{graphicx}
\usepackage{amsmath, amssymb}
\geometry{margin=1in}
\begin{document}
\begin{center}
History - K-12 Level Assessment 13 \
Total Marks: 40 \
Answer all questions.
\end{center}

\section*{Multiple Choice (10 marks)}
\begin{enumerate}[label=\arabic*.]
\item Homo erectus is most closely associated with which tool-making tradition?
\begin{enumerate}[label=(\alph*)]
    \item Occipital
    \item Acheulean
    \item Oldowan
    \item Stadial
\end{enumerate}
\item Which structures housed religious ceremonies?
\begin{enumerate}[label=(\alph*)]
    \item pit-houses
    \item pueblos
    \item kivas
    \item courtyard groups
\end{enumerate}
\item What was NOT a deciding factor in the development of Mesopotamian civilization?
\begin{enumerate}[label=(\alph*)]
    \item the need to concentrate population away from arable lands near rivers
    \item the need to construct monumental works, such as step-pyramids and temples
    \item the need to develop a complex social system that would allow the construction of canals
    \item the use of irrigation to produce enough food
\end{enumerate}
\item Genetic analysis shows that humans and chimps have been evolving separately for about how long?
\begin{enumerate}[label=(\alph*)]
    \item 2 million years
    \item 4 million years
    \item 7 million years
    \item 15 million years
\end{enumerate}
\item Which subfield of anthropology aims to better understand the nearest living relatives of humans?
\begin{enumerate}[label=(\alph*)]
    \item paleoanthropology
    \item archaeology
    \item primatology
    \item linguistics
\end{enumerate}
\end{enumerate}

\section*{True / False (10 marks)}
\textit{Answer True or False}
\begin{enumerate}[label=\arabic*.]
\setcounter{enumi}{5}
\item True or False: The Lapita pottery style is associated with a cultural complex that spread across the Pacific Islands.
\item True or False: Monk's Mound at Cahokia was primarily used as a burial monument for the leaders of the community.
\item True or False: Some archaeologists believe that Olmec monuments served primarily as surplus food storage centers for their communities.
\item True or False: The main Indus Valley cities had no defensive walls surrounding the upper cities.
\item True or False: Severe drought contributed to a decline in construction at Maya ceremonial centers during the civilization's history.
\end{enumerate}

\section*{Long Form Response (20 marks)}
\textit{Answer each question with a clear and complete explanation.}
\begin{enumerate}[label=\arabic*.]
\setcounter{enumi}{10}
\item Discuss the changes in Mayan civilization following the classical period around A.D. 800, focusing on its decline, relocation, resurgence, and subsequent decline.
\item Discuss the significance of radioactive dating techniques in determining the age of rocks, and explain how the half-life of isotopes like argon and potassium contributes to this process.
\end{enumerate}

\vfill
\noindent\textit{End of Paper}
\end{document}